\documentclass[11pt,a4paper]{article}
\usepackage[margin=1in]{geometry}
\usepackage{amsmath,amssymb,amsthm,graphicx,microtype}
\usepackage[hidelinks]{hyperref}
\newtheorem{theorem}{Theorem}
\newcommand{\R}{\mathbb R}
\title{The tangential $L\log L$ estimate fails at the printed $n=2$ endpoint}
\author{A counterexample to (4.1) in arXiv:1210.1454}
\date{17 August 2026}
\begin{document}
\maketitle

\noindent\textbf{Status.}
Candidate counterexample, likely valid, pending human review.  Estimate (4.1),
printed without a restriction beyond $n\geq2$, is false already for $n=2$,
even on a compact set strictly inside a smooth domain and for a sequence
converging strongly to zero in $W^{1,1}$.  The case $n\geq3$ is not settled.

\section{Source question}
Ka\l amajska, Kr\"omer, and Kružík ask whether the nonnegative tangential Jacobian
\[
 \det{}'\xi=\det(\xi_{ij})_{i,j=1,\ldots,n-1},
 \qquad \xi\in\R^{(n-1)\times n},
\]
obeys the uniform estimate
\begin{equation}\label{eq:source}
 \sup_k\int_K\gamma(\det{}'\nabla u_k)\,dx
 \leq C\bigl(K,\|u_k\|_{W^{1,n-1}(\Omega;\R^{n-1})}\bigr),
 \qquad \gamma(s)=s\log^+s,
\end{equation}
for compact $K\subset\overline\Omega$ away from boundary points whose outer
normal is not $e_n$ \cite[Section 4, equation (4.1), p.~12]{Source}.  The source
calls this an open problem and does not exclude $n=2$; its following
Example~4.1 explicitly starts with $n\geq2$.

\begin{center}
\includegraphics[width=\textwidth]{figures/open_problem_crop.png}
\end{center}

For a uniform estimate, the right side is naturally read as depending on a
uniform Sobolev bound.  The construction also defeats the per-index reading:
its Sobolev norms tend to zero while the integrals tend to infinity.

\section{Proof intuition}
At $n=2$, the matrix in $\det{}'$ is $1\times1$:
\[
 \det{}'\nabla u=\partial_1u,\qquad u:\Omega\to\R.
\]
No multidimensional determinant compensation remains.  Put a positive slope
$A_k$ on an interval of width $\ell_k$.  Its $W^{1,1}$ cost is
$A_k\ell_k$, whereas its $L\log L$ cost is
$A_k\ell_k\log A_k$.  Taking $A_k=e^{k^2}$ and
$A_k\ell_k=1/k$ makes the first vanish and the second equal $k$.

\section{Counterexample}
\begin{theorem}\label{thm:main}
Let $\Omega=B_2(0)\subset\R^2$ and
\[
 K=[-1/4,1/4]\times[-1/2,1/2]\Subset\Omega.
\]
There are $u_k\in W^{1,1}(\Omega)$, $k\geq2$, such that
\[
 u_k\to0\ \hbox{strongly in }W^{1,1}(\Omega),\qquad
 \det{}'\nabla u_k\geq0\ \hbox{a.e.},
\]
but
\[
 \int_K\gamma(\det{}'\nabla u_k)\,dx=k\to\infty.
\]
Consequently, the unqualified estimate~\eqref{eq:source} is false.
\end{theorem}

\begin{proof}
Set
\[
 A_k=e^{k^2},\qquad \ell_k=\frac1{kA_k}=\frac{e^{-k^2}}k,
\]
and define $u_k(x_1,x_2)=\psi_k(x_1)$, where
\[
 \psi_k(t)=
 \begin{cases}
 0,&t\leq0,\\
 A_kt,&0<t<\ell_k,\\
 1/k,&t\geq\ell_k.
 \end{cases}
\]
Almost everywhere,
\[
 \partial_1u_k=A_k\mathbf1_{(0,\ell_k)}(x_1),\qquad
 \partial_2u_k=0.
\]
Thus $\det{}'\nabla u_k=\partial_1u_k\geq0$.

Since $0\leq u_k\leq1/k$ and $|\Omega|=4\pi$,
\[
 \|u_k\|_{L^1(\Omega)}\leq\frac{4\pi}{k}.
\]
Every vertical section of $B_2(0)$ has length at most $4$, whence
\[
 \|\nabla u_k\|_{L^1(\Omega)}
 \leq4A_k\ell_k=\frac4k.
\]
Therefore $\|u_k\|_{W^{1,1}}\leq(4\pi+4)/k\to0$.

For $k\geq2$, $\ell_k<1/4$, so
$(0,\ell_k)\times[-1/2,1/2]\subset K$.  Since $\log A_k=k^2$,
\[
 \int_K\gamma(\det{}'\nabla u_k)\,dx
 =\ell_kA_k\log A_k
 =\frac1{kA_k}A_kk^2=k.
\]
Finally $K\Subset\Omega$, so it has positive distance from the whole boundary
and automatically meets the source condition on boundary normals.
\end{proof}

This is stronger than a bounded-sequence counterexample: adding weak
convergence to the question would not repair the endpoint.

\section{Scope and upgrade attempts}
The construction exploits the collapse of $\det{}'$ to a first derivative at
$n=2$.  It neither contradicts Mueller's theorem for genuinely
multidimensional nonnegative Jacobians nor disproves~\eqref{eq:source} for
$n\geq3$.

Five higher-dimensional routes were checked.  The source anisotropic example
has a nonintegrable normal derivative.  Truncating that singularity makes the
normal-derivative cost suppress the logarithmic divergence.  Reflection across
a flat boundary leaves a codimension-one minor, not the full Jacobian covered
by Mueller's theorem.  A diagonal embedding of the scalar spike forces the
$L^{n-1}$ energy to dominate its determinant mass.  Separated-variable and
shrinking-bubble scalings encounter the same tradeoff.  No proof or
counterexample resulted for $n\geq3$.

The final upgrade strengthened the low-dimensional example: choosing
$A_k\ell_k=1/k$ rather than a fixed mass makes $u_k\to0$ strongly while
$A_k\ell_k\log A_k\to\infty$.  This rules out weak-convergence and boundedness
repairs of the printed statement.

\section{Novelty and verification}
A bounded primary-source search used the exact source title, equation (4.1),
the phrases higher integrability and null Lagrangians at the boundary,
and the terms tangential Jacobian and $L\log L$.  Twelve indexed later
citations through 2020 were screened by title and primary text where available;
none explicitly resolves (4.1).  No later primary statement of this $n=2$
observation was found.  Novelty confidence is moderate-low because the example
is elementary; correctness confidence is high.

The included verifier checks the exact scaling formulas.  Human review should
first decide whether the source intended its printed $n\geq2$ range.  If it
silently intended $n\geq3$, this packet is a scope correction, not a solution
of that higher-dimensional problem.

\begin{thebibliography}{1}
\bibitem{Source}
A.~Ka\l amajska, S.~Kr\"omer, and M.~Kružík,
\emph{Sequential weak continuity of null Lagrangians at the boundary},
Calc. Var. Partial Differential Equations 49 (2014), 1263--1278;
arXiv:1210.1454.
\end{thebibliography}
\end{document}
